\subsection{\texorpdfstring{第~\ref{sec:completely-integrable-systems}~节}{完全可积系统}}\label{subsec:solutions-completely-integrable-systems}

\begin{description}[font=\normalfont,leftmargin=2em,labelwidth=1.5em]
\RGsolution{ex:prove-involutive-integral-dimension-bound}{involutive-integral-dimension-bound}
设 $F_1,\ldots,F_m\in C^\infty(T^*M)$ 对合, 并且在某点 $\alpha\in T^*M$ 处独立, 设 $X_{F_1},\ldots,X_{F_m}$ 是相应的 Hamilton 向量场. 则在 $\alpha$ 处, 向量 $(X_{F_i})_\alpha\in T_\alpha(T^*M)$ 线性无关. 实际上, 若对某些 $a_1,\ldots,a_m\in\mathbb{R}$ 有
\[
  \sum_{i=1}^m a_i(X_{F_i})_\alpha=0,
\]
则
\[
  0=\iota\left(\sum_{i=1}^m a_i(X_{F_i})_\alpha\right)\omega_\alpha
     =\sum_{i=1}^m a_i\iota((X_{F_i})_\alpha)\omega_\alpha
     =-\sum_{i=1}^m a_i(\mathrm{d}F_i)_\alpha,
\]
所以 $a_1=\cdots=a_m=0$, 因为 $(\mathrm{d}F_1)_\alpha,\ldots,(\mathrm{d}F_m)_\alpha$ 线性无关. 另一方面,
\begin{equation}\label{eq:involutive-hamiltonian-fields-are-isotropic}
  \omega_\alpha((X_{F_i})_\alpha,(X_{F_j})_\alpha)=\{F_i,F_j\}(\alpha)=0.
\end{equation}
取 $m$ 维空间 $W=\operatorname{span}\{(X_{F_1})_\alpha,\ldots,(X_{F_m})_\alpha\}$ 及其辛正交空间
\[
  W^\omega=\{w\in T_\alpha(T^*M)\mid
             \omega_\alpha(v,w)=0\ \text{对所有 }v\in W\}.
\]
则
\[
  2n=\dim T_\alpha(T^*M)=\dim W+\dim W^\omega.
\]
实际上, 考虑映射
\[
  \begin{aligned}
  \Phi:T_\alpha(T^*M)&\to W^*\\
                  v&\mapsto(\iota(v)\omega_\alpha)|_W,
  \end{aligned}
\]
容易看出 $\ker\Phi=W^\omega$ 且 $\operatorname{im}\Phi=W^*$ (因为 $W^*$ 中任一元素都可以延拓为 $T_\alpha^*(T^*M)$ 中的元素, 而映射 $v\mapsto\iota(v)\omega_\alpha$ 是 $T_\alpha(T^*M)$ 与 $T_\alpha^*(T^*M)$ 之间的同构). 因此,
\[
  2n=\dim T_\alpha(T^*M)=\dim W^\omega+\dim W^*=\dim W^\omega+\dim W.
\]
另一方面, 由 \eqref{eq:involutive-hamiltonian-fields-are-isotropic} 显然有 $W\subset W^\omega$. 所以
\[
  2n=\dim W^\omega+\dim W\ge2\dim W=2m,
\]
因而 $n\ge m$.

\RGsolution{ex:complete-integrability-of-surface-of-revolution}{complete-integrability-of-surface-of-revolution}
\begin{enumerate}[label=(\alph*),ref=\alph*,leftmargin=*]
\item 已知 $M$ 的测地线是由 $L=\frac12\langle v,v\rangle$ 确定的作用量的临界点, 其中 $\langle\cdot,\cdot\rangle$ 是 $\mathbb{R}^3$ 的 Euclid 度量在 $M$ 上诱导的度量. 若 $i:M\to\mathbb{R}^3$ 是标准包含映射, 则
\begin{align*}
  i^*g&=i^*(\mathrm{d}x\otimes\mathrm{d}x+\mathrm{d}y\otimes\mathrm{d}y+\mathrm{d}z\otimes\mathrm{d}z)\\
      &=(f(z))^2\,\mathrm{d}\theta\otimes\mathrm{d}\theta
        +((f'(z))^2+1)\,\mathrm{d}z\otimes\mathrm{d}z,
\end{align*}
因为
\[
  x=f(z)\cos\theta\quad\text{且}\quad y=f(z)\sin\theta,
\]
所以
\[
  i^*\mathrm{d}x=\mathrm{d}(f(z)\cos\theta)
      =-f(z)\sin\theta\,\mathrm{d}\theta+f'(z)\cos\theta\,\mathrm{d}z
\]
且
\[
  i^*\mathrm{d}y=\mathrm{d}(f(z)\sin\theta)
      =f(z)\cos\theta\,\mathrm{d}\theta+f'(z)\sin\theta\,\mathrm{d}z.
\]
因此 $M$ 的测地线的 Lagrange 函数由动能给出:
\[
  L(\theta,z,v^\theta,v^z)
      =\frac12\bigl((f(z))^2(v^\theta)^2+((f'(z))^2+1)(v^z)^2\bigr).
\]
\item 由于
\begin{align*}
  \frac{\partial L}{\partial v^\theta}&=(f(z))^2v^\theta,&
  \frac{\partial L}{\partial v^z}&=((f'(z))^2+1)v^z,\\
  \frac{\partial L}{\partial\theta}&=0,&
  \frac{\partial L}{\partial z}&=f(z)f'(z)(v^\theta)^2+f'(z)f''(z)(v^z)^2,
\end{align*}
Euler–Lagrange 方程为
\[
  \begin{cases}
    f(z)\ddot\theta+2f'(z)\dot\theta\dot z=0,\\
    ((f'(z))^2+1)\ddot z+f'(z)f''(z)\dot z^2-f(z)f'(z)\dot\theta^2=0.
  \end{cases}
\]
若 $\dot\theta=0$, 第一个 Euler–Lagrange 方程显然满足. 此外, 第二个方程化为
\[
  \frac{\ddot z}{\dot z}+\frac{f'(z)f''(z)}{(f'(z))^2+1}\dot z=0,
\]
进而
\[
  \frac{\mathrm{d}}{\mathrm{d}t}\left(\log\dot z+\frac12\log((f'(z))^2+1)\right)=0.
\]
所以对某个正的积分常数 $k$, 有
\[
  ((f'(z))^2+1)^{1/2}\dot z=k,
\]
从而
\[
  \frac{\mathrm{d}}{\mathrm{d}t}\left(\int_{z_0}^z((f'(s))^2+1)^{1/2}\,\mathrm{d}s\right)=k.
\]
注意
\[
  G(z)=\int_{z_0}^z((f'(s))^2+1)^{1/2}\,\mathrm{d}s
\]
是 $z$ 的递增函数 (因为 $((f'(s))^2+1)^{1/2}$ 为正), 因而是单射. 所以得到在局部坐标中由
\[
  \begin{cases}
    \theta(t)=\theta_0,\\
    z(t)=G^{-1}(kt)
  \end{cases}
\]
给出的轨迹, 其中 $\theta_0\in\mathbb{R}$ 为某一常数.

另一方面, 若轨迹满足 $f'(z(t))=0$, 则 Euler–Lagrange 方程变成
\[
  \ddot\theta=0\quad\text{且}\quad\ddot z=0.
\]
故对某些积分常数 $k,l$, 有 $\theta=\theta_0+kt$ 和 $z(t)=z_0+lt$. 因为要求 $f'(z(t))=0$, 若 $f$ 不是常函数, 则得到 $l=0$ 和 $z(t)=z_0$, 其中 $z_0$ 是 $f$ 的临界点, 并得到轨迹
\[
  \theta=\theta_0+kt\quad\text{且}\quad z(t)=z_0.
\]
若 $f$ 是常函数 (即 $S$ 是圆柱面), 则满足 $\theta=\theta_0+kt$ 和 $z(t)=z_0+lt$ 的任意轨迹都是解.
\item 在这些坐标下, Legendre 变换由
\[
  \begin{cases}
    p_\theta=\dfrac{\partial L}{\partial v^\theta}=(f(z))^2v^\theta,\\[4pt]
    p_z=\dfrac{\partial L}{\partial v^z}=((f'(z))^2+1)v^z
  \end{cases}
  \iff
  \begin{cases}
    v^\theta=\dfrac{p_\theta}{(f(z))^2},\\[4pt]
    v^z=\dfrac{p_z}{(f'(z))^2+1}
  \end{cases}
\]
给出. 因为它显然可逆, 所以 $L$ 是超正则的. 此时 Hamilton 函数为
\begin{align*}
  H(\theta,z,p_\theta,p_z)&=p_\theta v^\theta+p_zv^z-L
    =\frac{p_\theta^2}{(f(z))^2}+\frac{p_z^2}{(f'(z))^2+1}-L\\
  &=\frac{p_\theta^2}{2(f(z))^2}+\frac{p_z^2}{2((f'(z))^2+1)}.
\end{align*}
\item 由 Hamilton 方程,
\[
  \dot p_\theta=-\frac{\partial H}{\partial\theta}=0,
\]
所以 $p_\theta$ 是第一积分. 现在
\begin{align*}
  \mathrm{d}H
  &=\frac{p_\theta}{(f(z))^2}\,\mathrm{d}p_\theta
    +\frac{p_z}{(f'(z))^2+1}\,\mathrm{d}p_z\\
  &\quad-\left[\frac{f'(z)p_\theta^2}{(f(z))^3}
      +\frac{f'(z)f''(z)p_z^2}{((f'(z))^2+1)^2}\right]\mathrm{d}z,
\end{align*}
因而 $\mathrm{d}H$ 和 $\mathrm{d}p_\theta$ 在 $T^*M$ 的一个稠密开集上线性无关, 该集由坐标 $(z,\theta)$ 有定义且不满足
\[
  p_z=f'(z)\left[\frac{p_\theta^2}{(f(z))^3}
                +\frac{f''(z)p_z^2}{((f'(z))^2+1)^2}\right]=0
\]
的点组成. 例如, 只要 $p_z\ne0$, 即在某个 3 维子流形之外, 它们就线性无关. 所以 $H$ 和 $p_\theta$ 在稠密开集上独立.
\item 这个水平集的方程为 $p_\theta=l$ 和
\begin{align*}
  H=E&\iff\frac{l^2}{(f(z))^2}+\frac{p_z^2}{(f'(z))^2+1}=2E\\
     &\iff\frac{p_z^2}{(f'(z))^2+1}=2E-\frac{l^2}{(f(z))^2}.
\end{align*}
在使右端非负的点集上, 这些方程可解出 $p_z$, 即对于
\[
  f(z)\ge\frac l{\sqrt{2E}}.
\]
若 $f$ 在 $z=z_0$ 处有严格局部极大值, 则在像为 $z=z_0$ 的测地线附近, 不变水平集 $L_{(E,l)}$ 的投影形如 $z_{\min}\le z\le z_{\max}$, 其中 $z_{\min},z_{\max}$ 都接近 $z_0$, 且 $z_{\min}<z_0<z_{\max}$. 因此, 若测地线的初始条件接近 $z=z_0$ 的某个切向量, 则它将一直靠近 $z=z_0$, 这意味着该测地线是稳定的.
\end{enumerate}
\end{description}
